\documentclass[11pt]{article}
\usepackage[margin=1.1in]{geometry}
\usepackage{amsmath,amssymb,amsthm,mathtools}
\usepackage{url}
\usepackage[hidelinks]{hyperref}
\usepackage{microtype}

\newtheorem{theorem}{Theorem}
\newtheorem{proposition}[theorem]{Proposition}
\newtheorem{lemma}[theorem]{Lemma}
\newtheorem{corollary}[theorem]{Corollary}
\theoremstyle{definition}
\newtheorem{definition}[theorem]{Definition}
\theoremstyle{remark}
\newtheorem{remark}[theorem]{Remark}

\newcommand{\R}{\mathbb{R}}
\newcommand{\Q}{\mathbb{Q}}
\newcommand{\B}{B}
\DeclareMathOperator{\m}{m}
\DeclareMathOperator{\adj}{adj}
\DeclareMathOperator{\spn}{span}

\title{Six concyclic points satisfy the Euclidean Ramsey tensor criterion}
\author{JD Jones}
\date{October 2026}

\begin{document}
\maketitle

\begin{abstract}
A finite set of points in Euclidean space is Ramsey if every finite colouring of a sufficiently
high-dimensional Euclidean space contains a monochromatic congruent copy of it. The classification
of finite Euclidean Ramsey configurations announced in the openai/math release of 2026 characterises
the Ramsey sets that affinely span their space by a tensor criterion over $F \otimes_\Q F$, $F$ the
field generated by the coordinates, and proves that every set of at most five points on a circle is
Ramsey; a seven-point circle set is known not to be. We show that every set of six distinct
concyclic points satisfies the criterion, by an explicit certificate, and more generally every
spherical set of at most six points that spans its space and every spherical set of at most five
points. We also show that every spherical set with algebraic coordinates satisfies the criterion,
and that the criterion always implies sphericity. Given the classification, every spherical set of
at most six points is Ramsey, seven is the least size of a non-Ramsey spherical set, and the
spherical conjecture of Erd\H{o}s, Graham, Montgomery, Rothschild, Spencer and Straus holds for
spanning configurations with algebraic coordinates (and, through an algebraic spanning congruent copy,
for all of them). The unconditional statements are proved in Lean~4, and the Ramsey statements with the
classification as an explicit hypothesis; the heptagon's failure of the criterion and the corollary
that seven is least are on paper, the former with an exact checker.
\end{abstract}

\section{The criterion}

Let $A = \{a_1, \dots, a_s\} \subset \R^d$ be a finite set of distinct points, $F = \Q(a_{ij}) \subset \R$
the field generated by their coordinates, $\B = F \otimes_\Q F$, and $\m \colon \B \to F$ the
multiplication map $x \otimes y \mapsto xy$. Write $p_i = (1, a_i) \in F^{d+1}$ and index
$(d+1) \times (d+1)$ matrices by $\{0, 1, \dots, d\}$, the index $0$ being the constant coordinate.

\begin{definition}[the tensor criterion]\label{def:crit}
$A$ satisfies the criterion if there is $P \in \mathrm{Mat}_{d+1}(\B)$ with
\[
\sum_{\alpha,\beta} (p_{i\alpha} \otimes 1)\, P_{\alpha\beta}\, (1 \otimes p_{i\beta}) = 0
\quad\text{for every } i,
\qquad
\m(P_{\alpha\beta}) = \delta_{\alpha\beta} \quad\text{for } \alpha, \beta \geq 1 .
\]
\end{definition}

The constant row and column of $\m(P)$ are unconstrained. Applying $\m$ to the first identity gives
$p_i^{\mathsf T} \m(P) p_i = 0$: the points lie on the $F$-rational quadric $|x|^2 - 2\langle c', x\rangle + C = 0$
read off from $\m(P)$, which is a sphere. So the criterion asks that an $F$-rational sphere through
the points lift to $\B$ with the evaluation identities holding before multiplication.

The classification paper \cite{oai} states (Theorem~1.1 there) that a set of at least two points that
affinely spans $\R^d$ is Ramsey if and only if it satisfies the criterion, and proves (Corollary~7.4)
that every set of at most five points on a circle is Ramsey: for at most five concyclic points the
quadratic evaluation rows $p_i p_i^{\mathsf T}$ are linearly independent over $F$ (a product of two
chords isolates each point), and independent rows lift by an adjugate (Proposition~7.3). Six
concyclic points are the first case this does not reach: six points of a conic have rows spanning a
space of dimension five. The heptagon paper \cite{palvolgyi} gives seven points on a circle that are
not Ramsey. We use the classification only as a hypothesis in Section~\ref{sec:ramsey}.

In Lean the definitions are those of the classification's challenge file, restated character for
character; the repository is \texttt{six-concyclic-lean}, and the Lean name of each statement is given
in brackets.

\section{Six points on a circle}\label{sec:six}

\begin{theorem}[\texttt{six\_concyclic\_fieldCriterion}]\label{thm:six}
Every set of six distinct points on a circle in $\R^2$ satisfies the criterion.
\end{theorem}

\begin{proof}
Let $a_0, \dots, a_5$ be the points. Since the real circle through three of them is unique and its
centre is a rational function of their coordinates, there is an $F$-rational circle
$|x|^2 - 2\langle c', x\rangle + C = 0$ through all six (in Lean this comes from an $F$-linear
retraction $\R \to F$ fixing $1$, applied to the real equation; \texttt{exists\_rational\_sphere}).
Let $H$ be its matrix: spatial block $I_2$, constant row and column $-c'$, corner $C$.

\emph{Parametrisation.} Take the base point $p = a_5$ and $e = p - c'$, which is nonzero because the
points are distinct (\texttt{chord\_dot\_radius\_ne\_zero}). For a direction
$x = (u, w) \in F^2 \setminus 0$ the second intersection of the line $p + tx$ with the circle is
$q = p - 2\,\frac{\langle e, x\rangle}{|x|^2}\, x$, and homogeneously
\[
|x|^2 \, (1, q) = \Lambda \, (u^2, uw, w^2)^{\mathsf T},
\qquad
\Lambda =
\begin{pmatrix} 1 & 0 & 1 \\ p_0 - 2e_0 & -2e_1 & p_0 \\ p_1 & -2e_0 & p_1 - 2e_1 \end{pmatrix},
\]
with $\det \Lambda = 4|e|^2 \neq 0$ (\texttt{Lam\_ver}, \texttt{Winv\_mul\_Lam}). The tangent direction
$x = (-e_1, e_0)$ gives $|e|^2 (1, p)$ (\texttt{Lam\_tangent}). Put $x_j = a_j - p$ for $j < 5$ and
$x_5 = (-e_1, e_0)$, so that $p_j = N_j^{-1} \Lambda v(x_j)$ with $N_j = |x_j|^2 \ne 0$ and
$v(u,w) = (u^2, uw, w^2)$. Since $v(x)^{\mathsf T} \Lambda^{\mathsf T} H \Lambda v(x) = 0$ for every $x$
and the symmetric form $\Lambda^{\mathsf T} H \Lambda$ vanishes on the Veronese curve, it is a multiple of the
discriminant form $\delta(X, Y)^2$, $\delta = X_u Y_w - X_w Y_u$, whose matrix in the Veronese basis is
$K = \begin{psmallmatrix} 0 & 0 & 1 \\ 0 & -2 & 0 \\ 1 & 0 & 0\end{psmallmatrix}$; explicitly
$\Lambda^{\mathsf T} H \Lambda = -2|e|^2 K$ (\texttt{pullback}).

\emph{The bilinear forms.} For three directions $x_j$, $j \in S$, let
\[
\ell_S(X, Y) = \det
\begin{pmatrix}
X_u Y_u & X_u Y_w & X_w Y_u & X_w Y_w \\
x_{j,u} \otimes x_{j,u} & x_{j,u} \otimes x_{j,w} & x_{j,w} \otimes x_{j,u} & x_{j,w} \otimes x_{j,w}
\end{pmatrix}_{j \in S}
\]
a $(1,1)$-form in $X = x \otimes 1$ and $Y = 1 \otimes x$ with coefficients in $\B$. It vanishes at
$(x_j \otimes 1, 1 \otimes x_j)$ for $j \in S$ because the first row then repeats a row
(\texttt{ell\_vanishes}; this holds in any commutative ring). Put
$q = \ell_{\{0,1,2\}} \, \ell_{\{3,4,5\}}$, a biquadratic form with a $3 \times 3$ Veronese matrix $Q$
over $\B$: $v(X)^{\mathsf T} Q\, v(Y) = q(X,Y)$ (\texttt{Qm\_form}). Then $q$ vanishes at all six points.
Under $\m$ the two copies coincide; the corner coefficients of each $\ell_S$ vanish and the middle
ones are $\pm V_S$, $V_S = \prod_{j<k \in S} \det(x_j, x_k)$ (\texttt{mult\_coeffs}), so
$\m(Q) = V_1 V_2\, K$. The brackets are nonzero: $\det(x_j, x_k)$ for two chords from $p$ is twice the
signed area of a triangle of three distinct points of a circle (\texttt{det\_ne\_zero}), and
$\det(x_j, x_5) = \langle x_j, e\rangle = -|x_j|^2/2$ (\texttt{chord\_dot\_radius\_ne\_zero}).

\emph{The certificate.} Let $W = \Lambda^{-1}$, $s = -2|e|^2/(V_1 V_2) \in F$, and
\[
P = (s \otimes 1)\, (W \otimes 1)^{\mathsf T}\, Q\, (1 \otimes W) .
\]
At the point $i$, $W p_i = N_i^{-1} v(x_i)$, so the evaluation is
$(s N_i^{-1} \otimes N_i^{-1})\, q(x_i \otimes 1, 1 \otimes x_i) = 0$ (\texttt{eval\_P}). And
$\m(P) = s V_1 V_2\, W^{\mathsf T} K W = W^{\mathsf T} (\Lambda^{\mathsf T} H \Lambda) W = H$
(\texttt{mult\_P}), whose spatial block is $I_2$.
\end{proof}

No genericity is needed: the only inputs are that the points are distinct and concyclic. The
certificate rests on universal polynomial identities in the left and right copies of the fourteen
parameters (the base point, the centre and the five directions); certificate~0001 in
\texttt{note/certificates} verifies these identities and twelve exact end-to-end evaluations
independently of Lean.

\section{Lifts through isolating quadrics}\label{sec:lift}

\begin{lemma}[\texttt{fieldCriterion\_of\_isolators}]\label{lem:lift}
Let $H$ be an $F$-rational sphere matrix through $A$ with spatial block $I$. If for every $i$ there
is an $F$-matrix $q_i$ with $p_j^{\mathsf T} q_i p_j = \delta_{ij}$ for all $j$, then $A$ satisfies the
criterion.
\end{lemma}

\begin{proof}
Let $D$ be the $s \times (d+1)^2$ matrix over $\B$ with rows $(p_{i\alpha} \otimes 1)(1 \otimes p_{i\beta})$,
$S$ the $(d+1)^2 \times s$ matrix with columns $q_i \otimes 1$, $N = DS$, so that $\m(N) = I$ by the
isolating property, $b = H \otimes 1$, and
$z = \det(N)\, b - S \adj(N) D b$. Then $Dz = \det(N) Db - N \adj(N) Db = 0$, and
$\m(z) = \det(\m N) H - \m(S) \adj(\m N) \m(D b) = H$ since $\m(Db) = (p_i^{\mathsf T} H p_i)_i = 0$.
Reshape $z$ into $P$.
\end{proof}

This is the argument of Proposition~7.3 of \cite{oai} without the choice of a minor. Isolating
quadrics are products of affine functionals. On a sphere, three distinct points are never collinear
(\texttt{gram\_ne\_zero}: if two chords from a common point were parallel, $w = \lambda u$, the sphere
equation in the form $|x-p|^2 + 2\langle p - c', x - p\rangle = 0$ at both endpoints gives
$\lambda(\lambda-1)|u|^2 = 0$), so the line through two of the points has an $F$-rational affine
equation that is nonzero at any third (\texttt{exists\_line\_functional}), and a point is separated
from any other by $\langle a_i - a_j, x - a_j\rangle$.

\begin{theorem}[\texttt{at\_most\_five\_spherical\_fieldCriterion}]\label{thm:five}
Every spherical set of at most five distinct points, in any dimension, satisfies the criterion.
\end{theorem}

\begin{proof}
For each $i$ split the other at most four points into two groups of at most two; the product of
the two group functionals (a line, a point, or the constant $1$) vanishes at the others and not at
$a_i$. Apply Lemma~\ref{lem:lift} with the sphere of \texttt{exists\_rational\_sphere}.
\end{proof}

\begin{theorem}[\texttt{at\_most\_six\_spherical\_spanning\_fieldCriterion}]\label{thm:sixspan}
Every spherical set of at most six distinct points that affinely spans $\R^d$ satisfies the criterion.
\end{theorem}

\begin{proof}
At most five points: Theorem~\ref{thm:five}. Six points: $d \leq 1$ is impossible, $d = 2$ is
Theorem~\ref{thm:six}. For $d \geq 3$ and each $i$, some three of the other five points span a plane
not containing $a_i$ (\texttt{exists\_isolating\_triple}): fix two of them, $a_j, a_k$; if the plane
$\Pi$ through $a_j, a_k$ and a third point contains $a_i$, then any further point $a_m \notin \Pi$ gives
a plane through $a_j, a_k, a_m$ meeting $\Pi$ in the line $a_j a_k$, which does not contain $a_i$; and
if every point lies in $\Pi$ the set does not span. The affine function vanishing on that plane
(the component of $a_i - a_j$ orthogonal to the plane, computed over $F$ with the $2 \times 2$ Gram
matrix of two chords; \texttt{exists\_plane\_functional}) times the line through the remaining two
points isolates $a_i$; apply Lemma~\ref{lem:lift}.
\end{proof}

\begin{corollary}[\texttt{at\_most\_six\_concyclic\_fieldCriterion}]
Every set of at most six distinct points on a circle satisfies the criterion.
\end{corollary}

\section{Algebraic coordinates}\label{sec:alg}

\begin{theorem}[\texttt{algebraic\_spherical\_fieldCriterion}]\label{thm:alg}
Every spherical set whose coordinates are algebraic numbers satisfies the criterion, whatever its size.
\end{theorem}

\begin{proof}
$F$ is a number field, so $\B = F \otimes_\Q F$ is a finite \'etale $F$-algebra and the multiplication
map splits: there is an idempotent $t \in \B$ with $(1 \otimes x - x \otimes 1)\, t = 0$ for every $x \in F$
and $\m(t) = 1$ (the separability idempotent; in Mathlib,
\texttt{Algebra.FormallyUnramified.iff\_exists\_tensorProduct}). For an $F$-rational sphere matrix $H$
put $P = t \cdot (H \otimes 1)$. In the sandwich, $t$ moves every right factor to the left:
$(x \otimes 1)\, t\, (H_{\alpha\beta} \otimes 1)\, (1 \otimes y) = t \cdot (x H_{\alpha\beta} y \otimes 1)$, so the
evaluation at $a_i$ is $t \cdot (p_i^{\mathsf T} H p_i \otimes 1) = 0$; and $\m(P) = H$.
\end{proof}

\begin{theorem}[\texttt{fieldCriterion\_spherical}, \texttt{algebraic\_fieldCriterion\_iff\_spherical}]
The criterion implies sphericity; for algebraic coordinates the criterion is exactly sphericity.
\end{theorem}

\begin{proof}
The first is the remark after Definition~\ref{def:crit}, with the $F$-rational quadric read in $\R$:
$|x - c'|^2 = |c'|^2 - C$, a sphere since a real point satisfies it. The second adds
Theorem~\ref{thm:alg}.
\end{proof}

So the content of Theorems~\ref{thm:six}--\ref{thm:sixspan} lies in the transcendental case: for a
coordinate field that is not algebraic over $\Q$ no separability idempotent exists (a field that is a
separable $\Q$-algebra is finite over $\Q$), though $\B$ may still have other idempotents, as for
$F = \Q(t, \sqrt 2)$; when $\Q$ is algebraically closed in $F$, $\B$ is a domain. Every known non-Ramsey
spherical set has a transcendental coordinate, as it must if (S) below holds.

\section{Consequences given the classification}\label{sec:ramsey}

Write (S) for the sufficiency half of the classification: a configuration of at least two distinct
points that affinely spans $\R^d$, $d \geq 1$, and satisfies the criterion is Ramsey. The Ramsey
property passes to subsets and to congruent configurations, and any configuration of at least two
distinct spherical points is congruent to one that spans a space of positive dimension and is still
spherical (translate, take the span, map it isometrically; the centre projects orthogonally onto the
span; \texttt{exists\_spanning\_representative}).

\begin{theorem}[\texttt{at\_most\_six\_spherical\_ramsey\_of\_sufficiency}]\label{thm:ramsey6}
Given \emph{(S)}, every spherical set of at most six points is Ramsey.
\end{theorem}

\begin{proof}
Theorem~\ref{thm:sixspan} on the spanning representative, then (S) and congruence.
\end{proof}

\begin{theorem}[\texttt{algebraic\_spherical\_ramsey\_of\_sufficiency}]
Given \emph{(S)}, every spanning spherical set with algebraic coordinates is Ramsey; given both halves of the
classification, a spanning set with algebraic coordinates is Ramsey if and only if it is spherical
(the Lean name of the second statement is \texttt{algebraic\_ramsey\_iff\_}\allowbreak
\texttt{spherical\_of\_classification}).
\end{theorem}

The first statement is the spherical conjecture (asked in \cite{egmrss}, stated as a conjecture in
\cite{graham}) for spanning configurations with algebraic coordinates, given (S); a non-spanning
configuration with algebraic coordinates has a spanning congruent copy with algebraic coordinates
(Gram--Schmidt over the real algebraic numbers), which is not formalised, so the conjecture follows for
all algebraic configurations on paper. With the theorem of \cite{egmrss} that Ramsey sets are spherical,
the second statement is the equivalence. Neither \cite{oai} nor its scope note states either.

\begin{proposition}\label{prop:hept}
The seven-point set $\{(0,0)\} \cup \{(j, \pm\sqrt{j(2r-j)}) : j \in \{1,3,4\}\}$ of \cite{palvolgyi}, $r > 2$
transcendental, does not satisfy the criterion.
\end{proposition}

\begin{proof}
Its coordinate field is $F = \Q(r, y_1, y_3, y_4)$ with $y_j^2 = j(2r-j)$; let $D$ be the derivation of
$F$ with $D(r) = 1$, so $D(y_j) = j/y_j$, and $\Phi(x \otimes y) = D(x) D(y)$, which is well defined on
$\B$ since it is $\Q$-bilinear. Give the origin weight $2$ and the pairs $j = 1, 3, 4$ weights $-2, 2, -1$.
Then, exactly,
\[
\textstyle\sum \lambda_i\, \tilde p_i \tilde p_i^{\mathsf T} = 0, \qquad
\sum \lambda_i\, D(\tilde p_i)\, \tilde p_i^{\mathsf T} = 0, \qquad
\sum \lambda_i\, |D \tilde p_i|^2 = \frac{-24}{(2r-1)(2r-3)(2r-4)} \neq 0,
\]
with $\tilde p_i = (1, a_i)$ (certificate 0002). Besides $\Phi$ and $\m$, the $\Q$-bilinear maps
$\Psi_1(u \otimes v) = D(u)\, v$ and $\Psi_2(u \otimes v) = u\, D(v)$ are well defined on $\B$, and for
every $z \in \B$ the Leibniz rule gives
$\Phi\big((x \otimes 1) z (1 \otimes y)\big) = D(x) D(y)\, \m(z) + D(x)\, y\, \Psi_2(z) + x\, D(y)\, \Psi_1(z) + xy\, \Phi(z)$
(check it on pure tensors and extend by linearity). For any $P$ satisfying the evaluation identities, sum
this against $\lambda_i$ over the points and over $\alpha, \beta$ with $x = \tilde p_{i\alpha}$,
$y = \tilde p_{i\beta}$, $z = P_{\alpha\beta}$: the first moment identity kills the $\Phi$ term, the second
and its transpose kill the two mixed terms, and what remains is $\sum_{\alpha\beta} \big(\sum_i \lambda_i D(\tilde p_{i\alpha}) D(\tilde p_{i\beta})\big) \m(P_{\alpha\beta})$.
With the spatial block of $\m(P)$ equal to $I$ (and $D(1) = 0$ on the constant row) this is
$\sum_i \lambda_i |D \tilde p_i|^2 \neq 0$, while the left side is $\Phi(0) = 0$.
\end{proof}

\begin{corollary}
Given (S), and either the necessity half of the classification or Theorem~1 of \cite{palvolgyi}, the
least size of a non-Ramsey spherical set is seven.
\end{corollary}

\begin{proof}
At most six: Theorem~\ref{thm:ramsey6}. The heptagon is not Ramsey by Proposition~\ref{prop:hept} and the
necessity half of the classification, or by Theorem~1 of \cite{palvolgyi} directly.
\end{proof}

\begin{remark}
The threshold is about the existence of a failing set, not about every seven-point set. Exact
computations made during the development, not included here, decided the criterion over the field
$\Q(s)$ (where $F \otimes_\Q F$ localised at the kernel of $\m$ is a discrete valuation ring) and found
seven- and eight-point sets on the unit circle with stereographic parameters $c_i s + d_i$, $c_i, d_i$
rational, that satisfy the criterion, while the nine-point and twelve-point sets of \cite{oai} fail it.
The heptagon's failure rests on an algebraic relation among its coordinates; that generic seven points
also fail is stated without proof in the appendix of \cite{palvolgyi}.
\end{remark}

\section{Prior art and provenance}

As of 2026-10-07 (UTC) no statement that six concyclic points are Ramsey or satisfy the criterion, and no
statement of the criterion for algebraic configurations, was found in \cite{oai} (its text, scope
note and Lean tree, which stop at five points), in \cite{palvolgyi} (whose appendix shows only that
derivation identities cannot refute at most six points), on arXiv, MathOverflow, erdosproblems.com
(problem 174), the formal-conjectures repository, the Palomar registry, Hexagon, or GitHub. The
classical positive results are the regular hexagon and all sets with a soluble transitive symmetry
group \cite{kriz1}, trapezoids \cite{kriz2}, simplices \cite{fr} and regular polytopes
\cite{cantwell}; the status of general cyclic quadrilaterals before 2026 and the subtransitivity
conjecture are in \cite{lrw}.

Claude (Anthropic) was used throughout; the theorems, proofs, searches and checkers arose in that
process, and the repository was written by Claude. The Lean development was checked by the Lean~4 kernel \cite{lean} against Mathlib
\cite{mathlib} at tag v4.35.0-rc2. The author directed the project and takes responsibility for the
claims. This note is licensed CC~BY-SA~4.0; the repository is MIT, with OpenAI's Apache-2.0 notice on
the restated definitions.

\begin{thebibliography}{9}
\bibitem{oai} OpenAI, \emph{A classification of finite Euclidean Ramsey configurations}, openai/math, family
172, preprint dated 23 September 2026, commit \texttt{adc7f124} (6 October 2026).
\bibitem{palvolgyi} D.~P\'alv\"olgyi, \emph{A cyclic non-Ramsey heptagon}, arXiv:2609.23327 (20 September 2026).
\bibitem{egmrss} P.~Erd\H{o}s, R.~L.~Graham, P.~Montgomery, B.~L.~Rothschild, J.~Spencer, E.~G.~Straus,
Euclidean Ramsey theorems.~I, J.~Combin.\ Theory Ser.~A 14 (1973) 341--363.
\bibitem{graham} R.~L.~Graham, Recent trends in Euclidean Ramsey theory, Discrete Math.\ 136 (1994) 119--127.
\bibitem{kriz1} I.~K\v{r}\'{\i}\v{z}, Permutation groups in Euclidean Ramsey theory, Proc.\ Amer.\ Math.\ Soc.\ 112 (1991) 899--907.
\bibitem{kriz2} I.~K\v{r}\'{\i}\v{z}, All trapezoids are Ramsey, Discrete Math.\ 108 (1992) 59--62.
\bibitem{fr} P.~Frankl, V.~R\"odl, A partition property of simplices in Euclidean space, J.~Amer.\ Math.\ Soc.\ 3 (1990) 1--7.
\bibitem{cantwell} K.~Cantwell, All regular polytopes are Ramsey, J.~Combin.\ Theory Ser.~A 114 (2007) 555--562.
\bibitem{lrw} I.~Leader, P.~A.~Russell, M.~Walters, Transitive sets and cyclic quadrilaterals, J.~Combin.\ 2 (2011) 457--462; Transitive sets in Euclidean Ramsey theory, J.~Combin.\ Theory Ser.~A 119 (2012) 382--396.
\bibitem{lean} L.~de~Moura, S.~Ullrich, The Lean 4 theorem prover and programming language, CADE 28, LNCS, Springer (2021) 625--635.
\bibitem{mathlib} The mathlib Community, The Lean mathematical library, CPP 2020, 367--381.
\end{thebibliography}

\end{document}
